%% ====================================================================
%%  Systemic Risk as Geometry: A Unified Theory of Financial Instability,
%%  Detection, and Adaptive Stabilisation
%%
%%  Draft v0.1  —  February 2026
%% ====================================================================
\documentclass[11pt]{article}

%% ---- Geometry -------------------------------------------------------
\usepackage[margin=1.15in, top=1.1in, bottom=1.2in]{geometry}

%% ---- Mathematics ---------------------------------------------------
\usepackage{amsmath, amsthm, amssymb, mathtools}
\usepackage{bm}          % bold math symbols
\usepackage{bbm}         % blackboard bold for 1

%% ---- References / Hyperlinks ---------------------------------------
\usepackage[colorlinks=true, linkcolor=blue!70!black,
            citecolor=blue!60!black, urlcolor=blue!60!black]{hyperref}
\usepackage[numbers, sort&compress]{natbib}

%% ---- Fonts / Layout ------------------------------------------------
\usepackage{microtype}
\usepackage{booktabs, tabularx, multirow}
\usepackage{graphicx}
\usepackage{xcolor}
\usepackage{tcolorbox}
\tcbuselibrary{skins, breakable}

%% ---- Algorithms ----------------------------------------------------
\usepackage{algorithm, algpseudocode}

%% ---- Appendix ------------------------------------------------------
\usepackage[toc, page]{appendix}

%% ====================================================================
%%  Theorem Environments
%% ====================================================================
\newtheorem{theorem}{Theorem}[section]
\newtheorem{lemma}[theorem]{Lemma}
\newtheorem{proposition}[theorem]{Proposition}
\newtheorem{corollary}[theorem]{Corollary}
\newtheorem{definition}[theorem]{Definition}
\newtheorem{remark}[theorem]{Remark}
\newtheorem{assumption}{Assumption}
\renewcommand{\theassumption}{E\arabic{assumption}}

\theoremstyle{definition}
\newtheorem{example}[theorem]{Example}

%% ====================================================================
%%  Macros
%% ====================================================================
\DeclareMathOperator{\sym}{sym}
\DeclareMathOperator{\erf}{erf}
\DeclareMathOperator*{\argmin}{arg\,min}
\newcommand{\R}{\mathbb{R}}
\newcommand{\E}{\mathbb{E}}
\newcommand{\Xst}{X^{\star}}
\newcommand{\grad}{\nabla}
\newcommand{\Ebs}{E_{\mathrm{BS}}}
\newcommand{\Cman}{\mathcal{C}}
\newcommand{\Lpair}{L_{\mathrm{pair}}}
\newcommand{\lmax}{\lambda_{\max}}
\newcommand{\lmin}{\lambda_{\min}}
\newcommand{\norm}[1]{\left\|#1\right\|}
\newcommand{\inner}[2]{\left\langle #1,\, #2\right\rangle}
\newcommand{\costheta}{\cos\theta}
\newcommand{\ddt}{\tfrac{d}{dt}}

%% ====================================================================
%%  Title / Authors
%% ====================================================================
\title{%
  \textbf{Systemic Risk as Geometry}\\[4pt]
  \large A Unified Theory of Financial Instability, Detection,
         and Adaptive Stabilisation
}
\author{%
  [Author redacted for review]
}
\date{Draft: \today \\ \small\textit{Subject to revision — not for circulation}}

%% ====================================================================
\begin{document}
\maketitle

\begin{abstract}
We establish a geometric equivalence between the statistical structure of financial
blind spots and the Lyapunov energy of networked agent dynamics.  The Blind-Spot
Detection Theory (BSDT) deviation operators, calibrated solely on data from normal
market conditions, generate a functional $\Ebs(X)$ whose gradient aligns with the
instability-restoring force of GravityEngine to within a factor of two
(empirical mean alignment angle $\costheta = 0.8616$; theoretical lower bound
$\costheta \geq 1 - \varepsilon$ under a separation condition on the operators).
We prove three consequences.
(i)~\textbf{Theorem~C} (equivalence): $\Ebs$ and the GravityEngine potential $\Phi$
are equivalent Lyapunov functions on every compact sublevel set; their gradients
are uniformly equivalent in magnitude.
(ii)~\textbf{Theorem~2} (phase transition): a spectral radius criterion identifies a
critical manifold $\Cman$ above which no bounded constant coupling stabilises the
system, but state-adaptive $\gamma(X)$ does.
(iii)~\textbf{Theorem~OR2} (optimal policy): the adaptive damping law
$\gamma^*(X) \approx E(X)/\bigl(E(X)+\theta\bigr)$ is the closed-form solution to an
infinite-horizon stochastic control problem with intervention cost $\kappa$.
Together these results justify state-contingent macroprudential policy as the
\emph{minimal} intervention above $\Cman$: weaker rules fail by the impossibility
theorem for constant coupling (Theorem~6); stronger rules incur excess intervention
cost bounded by the $O(\sqrt{T})$ regret bound (Proposition~CS1).
We calibrate the model to three crises—2008~GFC, 2020~COVID, 2023 regional banking
stress—and show $\gamma^*(X_t)$ leads standard risk measures (SRISK, $\Delta$CoVaR)
by 2--6 quarters.
\end{abstract}

\tableofcontents
\newpage

%% ====================================================================
\section{Introduction}
\label{sec:intro}
%% ====================================================================

\paragraph{The dual failure.}
Financial crises share a common signature: the institutions most exposed to systemic
risk are precisely the ones whose internal models fail to detect it.  In the 2008
Global Financial Crisis, major banks reported low Value-at-Risk in the quarters
immediately preceding collapse~\citep{brunnermeier2009deciphering}.  In the 2020
COVID shock, leverage in the U.S.\ banking system was at a post-crisis high
when the stress materialised~\citep{greenwald2020financial}.  Detection failure
and financial instability are not independent problems that happen to coincide---they
are the same geometric phenomenon viewed from two different angles.

This paper makes that claim precise.  We show that the statistical structure of
detection blind spots (formalised as the Blind-Spot Detection Theory deviation
operators, BSDT) and the Lyapunov energy of networked agent dynamics (GravityEngine)
are equivalent representations of the same object: the distance from normality in
the configuration space of financial institutions.  The central result,
Theorem~C, states that the gradient of the detection functional and the restoring
force of the physical model agree to within a constant factor---empirically
$\costheta = 0.862$, theoretical lower bound $\costheta \ge \nu/c_2$ under a
non-degeneracy condition on the operators.

\paragraph{Consequences.}
Three consequences follow.  First, a data-constructed Lyapunov function (calibrated
solely on normal-market observations) stabilises the physical dynamics: detection
and stabilisation are not separate engineering problems.  Second, there is a sharp
spectral phase transition at a critical manifold $\Cman$ above which no bounded
constant coupling can stabilise the system---this is the formal statement of
procyclicality, previously a qualitative critique of Basel~III.  Third, the adaptive
damping law $\gamma^*(X)$ that stabilises the system above $\Cman$ is the
closed-form solution to an infinite-horizon stochastic control problem, giving a
Pigouvian micro-foundation for state-contingent macroprudential policy.

\paragraph{Contributions.}
\begin{enumerate}
  \item \textbf{Theorem~C} (geometry): equivalence of BSDT gradient and GravityEngine
        restoring force, with explicit constants $c_1, c_2$.
  \item \textbf{Theorem~2} (phase transition): spectral radius criterion for the
        critical manifold; impossibility of constant coupling above it.
  \item \textbf{Proposition~CS1} (online algorithm): $O(\sqrt{T})$ regret for
        projected OGD on $\gamma_t$; $O(N^2 d \log N)$ per step.
  \item \textbf{Theorem~OR2} (optimal policy): closed-form $\gamma^*(X)$ from
        dynamic programming; Pigouvian interpretation.
  \item \textbf{Empirical validation} (Section~\ref{sec:empirical}): the model
        leads SRISK and $\Delta$CoVaR by 2--6 quarters across three crisis windows.
\end{enumerate}

\paragraph{Related work.}
The systemic risk measurement literature includes SRISK~\citep{brownlees2017srisk},
$\Delta$CoVaR~\citep{adrian2016covar}, and the CATFIN index~\citep{allen2012systemic}.
These are regression-based tail-risk measures; they are backward-looking (estimated
from historical returns) and do not provide stability guarantees.  Our approach is
forward-looking: the adaptive law $\gamma^*(X)$ is derived from the current
configuration $X_t$, not from a historical regression.  Network-theoretic approaches
(\citealt{acemoglu2015systemic}; \citealt{elliott2014financial}) characterise when
network structure amplifies shocks but do not provide a practical regulatory
operator.  The Lyapunov approach to financial stability is explored
in~\citealt{may2008complex} for ecological networks and adapting it to financial
systems with data-constructed potentials is, to our knowledge, new.
The OCO framing of regulatory policy is related to~\citealt{shalev2012online}
but the loss function here (energy change under coupling) differs from standard
online learning setups.

\begin{tcolorbox}[title=\textbf{Mathematical Overview}, colframe=blue!40!black,
                  colback=blue!5!white, breakable]
\textbf{Theorem~C} (Section~\ref{sec:equiv}): The BSDT blind-spot gradient and the
GravityEngine restoring force are equivalent in magnitude and direction.
Formally, $c_1\norm{\grad\Phi}\le\norm{\grad\Ebs}\le c_2\norm{\grad\Phi}$ and
$\costheta \ge 1 - \varepsilon$ on every compact sublevel set.

\smallskip
\textbf{Theorem~2} (Section~\ref{sec:damping}): A spectral phase transition at
$\lmax(\rho\ell\mathbf{W}) = 1$ separates the regime where constant damping
suffices from the regime where it fails; adaptive $\gamma(X)$ is necessary and
sufficient above $\Cman$.

\smallskip
\textbf{Proposition~CS1} (Section~\ref{sec:online}): Online gradient descent on
$\gamma_t$ achieves cumulative regret $R_T \le O(\sqrt{T})$ against the best fixed
policy in hindsight, with per-step cost $O(N^2 d \log N)$ — matching the existing
pairwise force computation.

\smallskip
\textbf{Theorem~OR2} (Section~\ref{sec:policy}): The infinite-horizon optimal
policy has the closed-form approximation
$\gamma^*(X) \approx \tfrac{\beta\norm{\grad E}^2}{2\kappa+\beta\norm{\grad E}^2}
\cdot \tfrac{E(X)}{E(X)+\theta}$,
pinning the free parameter $\theta$ to the ratio of intervention cost~$\kappa$
to discount factor~$\beta$.
\end{tcolorbox}

%% ====================================================================
\section{Agent Model and Equilibrium Dynamics}
\label{sec:model}
%% ====================================================================

We derive the force field $F(X)$ from the optimising behaviour of $N$ financial
institutions.  This section is the micro-foundation: it pins the abstract
dynamics to observable quantities (leverage, correlation, network topology).

\subsection{The $N$-Institution Game}

There are $N$ financial institutions (banks, funds, or insurance companies).
Each institution $i$ occupies position $x_i(t) \in \R^d$ in a $d$-dimensional
economic state space, where the coordinates represent standardised financial
characteristics (e.g.\ log-leverage, deposit-to-asset ratio, net interbank
exposure, return volatility).  The joint state is $X(t) \in \R^{N \times d}$.

\begin{definition}[Economic state space]
The coordinates of $x_i \in \R^d$ are chosen so that the \emph{normal operating
region} is an ellipsoid $\mathcal{N} = \{x : (x - \mu_0)^\top \Sigma_0^{-1}(x-\mu_0) \le r^2\}$
centred at the historical mean $\mu_0$ of the banking system.
The BSDT deviation operators (Section~\ref{sec:bsdt}) measure distance from $\mathcal{N}$.
\end{definition}

Each institution $i$ maximises expected profit subject to limited liability:
\begin{equation}\label{eq:institution_opt}
  \max_{\ell_i \ge 0}\; \E[\ell_i \cdot r_i(X)] - \tfrac{\mu}{2}\ell_i^2,
  \qquad r_i(X) = \bar r + \beta_i^\top X + \epsilon_i,
\end{equation}
where $\ell_i$ is leverage, $\mu > 0$ is risk-aversion, $\beta_i \in \R^d$ is the
factor loading of institution $i$ (its sensitivity to system state), and
$\epsilon_i \sim N(0,\sigma^2)$ is idiosyncratic noise.  The herding channel enters
through $\beta_i^\top X$: institutions share exposure to a common
system state.

\subsection{Nash Equilibrium and Positive Feedback}

The best-response leverage of institution $i$ given all others' choices is:
\[
  \ell_i^*(X) = \frac{\beta_i^\top X}{\mu} \cdot \mathbf{1}[\beta_i^\top X > 0].
\]
The resulting update to $x_i$ is proportional to the excess return:
$\dot x_i = \ell_i^*(X) \cdot [r_i(X) - \bar r]$, giving:
\begin{equation}\label{eq:micro_dynamics}
  \dot x_i = \frac{(\beta_i^\top X)^2}{\mu} \cdot \hat\beta_i
  \;+\; \text{(mean-reversion and pairwise terms)},
\end{equation}
where $\hat\beta_i = \beta_i/\norm{\beta_i}$.
The leading term is \emph{quadratic positive feedback}: when the system is displaced
from $\mu_0$ in the direction of $\beta_i$, institution $i$ \emph{amplifies} the
displacement.  This is the micro-foundation of the amplifying dynamics assumed in
the abstract model.

\subsection{The GravityEngine as Reduced-Form Dynamics}

Aggregating over $N$ institutions and adding mean-reversion and pairwise interaction
terms (modelling balance-sheet spillovers and correlated fire-sales), the system
dynamics become:
\begin{equation}\label{eq:full_dynamics}
  \dot X = \underbrace{-\alpha(X - \mu\mathbf{1}^\top)}_{\text{mean-reversion}}
           \underbrace{-\, \grad_X \Phi_{\mathrm{pair}}(X)}_{\text{balance-sheet spillovers}}
           \underbrace{-\, \sum_k \beta_k \cdot D\Phi_k(X)^\top D_k(X)}_{\text{operator-driven forces}}
           \;=: F(X).
\end{equation}
This is precisely the GravityEngine force decomposition ($F_{\mathrm{rad}} + F_{\mathrm{pair}} + F_{\mathrm{op}}$).
The parameter $\alpha$ is the speed of mean-reversion (quantifiable from BIS data as
the rate at which leverage ratios converge to steady-state after shocks).

\subsection{Critical Manifold from First Principles}

The Jacobian of $F$ at a configuration $X$ is $DF(X) = -\alpha I - D^2\Phi_{\mathrm{pair}}(X)$.
Define $\rho(X)$ as mean pairwise correlation and $\ell(X)$ as mean leverage.
In the mean-field limit ($N \to \infty$, homogeneous $\beta_i = \beta$):
\[
  \lmax(D^2\Phi_{\mathrm{pair}}) = \rho(X)\,\ell(X)\,\lmax(\mathbf{W})
\]
where $\mathbf{W}$ is the row-normalised adjacency matrix.  Instability occurs when
$\lmax(-\alpha I - D^2\Phi_{\mathrm{pair}}) > 0$, i.e.\ when
$\rho(X)\ell(X)\lmax(\mathbf{W}) > \alpha$.
Normalising to $\alpha = 1$ (WLOG by rescaling time) recovers the critical manifold:
\[
  \Cman = \bigl\{X : \lmax(\rho(X)\,\ell(X)\,\mathbf{W}) = 1\bigr\}.
\]
All three quantities $\rho, \ell, \mathbf{W}$ are directly observable: $\rho$ from
return correlations (CRSP), $\ell$ from balance sheets (FRED/Compustat), and
$\mathbf{W}$ from bilateral exposures (BIS).

\subsection{The Pigouvian Interpretation}

The decentralised equilibrium produces $E_{\mathrm{DE}}(X) > E_{\mathrm{SP}}(X)$: each
institution ignores the externality its leverage choice imposes on the system
energy $E(X)$.  The social planner's problem is to choose a tax $\tau_i(X)$ on
leveraged positions such that the corrected best-response recovers the socially
optimal dynamics.  The tax that achieves this is:
\[
  \tau_i^*(X) = \gamma^*(X) \cdot \frac{\partial E(X)}{\partial \ell_i}
\]
where $\gamma^*(X)$ is the adaptive coupling derived in Section~\ref{sec:policy}.
This is a Pigouvian correction: the tax equals the marginal externality.  The
optimal macroprudential instrument is therefore not a fixed capital surcharge but
a state-contingent levy proportional to the current energy gradient---precisely
the adaptive damping law.

%% ====================================================================
\section{Blind-Spot Energy Functional}
\label{sec:bsdt}
%% ====================================================================

The Blind-Spot Detection Theory (BSDT) formalises four modes by which standard
statistical detectors fail to identify abnormal financial behaviour.  We define each
operator, state its regularity, and construct the blind-spot energy functional.

\subsection{The Four Deviation Operators}

Fix a reference measure $\mathcal{N}$ (the distribution of agent states under normal
market conditions), with mean $\mu_0 \in \R^d$ and covariance $\Sigma_0 \in \R^{d \times d}$.
All four operators map agent configuration $X$ to non-negative scalars.

\begin{definition}[Camouflage deviation $\delta_C$]
An agent camouflages when it lies in a high-density region of $\mathcal{N}$ despite
having anomalous \emph{conditional} behaviour:
\[
  \delta_C(x_i) = \max\bigl(0,\; S(x_i) - \E_{\mathcal{N}}[S]\bigr),
\]
where $S(x_i) = (x_i - \mu_0)^\top \Sigma_0^{-1}(x_i - \mu_0)$ is the
Mahalanobis distance. $\delta_C$ is large when an agent appears normal by marginal
distance but its conditional residuals are large.
\end{definition}

\begin{definition}[Feature-gap deviation $\delta_G$]
The feature gap measures the share of variance in $x_i - \mu_0$ not explained by
the leading eigenvectors of $\Sigma_0$:
\[
  \delta_G(x_i) = \norm{x_i - \mu_0}^2 - \norm{\Pi_k(x_i - \mu_0)}^2,
\]
where $\Pi_k$ projects onto the top-$k$ principal components of $\mathcal{N}$.
$\delta_G$ is large when an agent moves in directions the normal distribution does
not cover---a structural anomaly invisible to PCA-based monitors.
\end{definition}

\begin{definition}[Activity-anomaly deviation $\delta_A$]
The activity operator measures excess trading volume or transaction frequency
relative to a normal baseline:
\[
  \delta_A(x_i) = \max\bigl(0,\; \norm{\dot x_i} - v_0\bigr),
\]
where $v_0$ is the 95th percentile of $\norm{\dot x}$ under $\mathcal{N}$.
An agent with $\delta_A > 0$ is moving faster than normal---often a precursor to
fire-sale dynamics.
\end{definition}

\begin{definition}[Temporal-novelty deviation $\delta_T$]
The temporal operator measures whether the current state $x_i(t)$ is a novel
visitation: low probability under a kernel density estimate built on the agent's
own history $\{x_i(s) : s < t\}$:
\[
  \delta_T(x_i, t) = \max\bigl(0,\; -\log p_h(x_i(t))\bigr),
\]
where $p_h$ is a Gaussian kernel density with bandwidth $h$.  $\delta_T$ is large
when an agent visits states it has never visited before---temporal novelty.
\end{definition}

\subsection{The Blind-Spot Energy Functional}

\begin{definition}[Blind-spot energy]
\label{def:ebs}
Let $\psi_i : \R^+ \to \R^+$ be a convex, non-decreasing, 1-Lipschitz link
function with $\psi_i(0) = 0$ (e.g.\ $\psi_i(s) = \tanh(s)$ or $\psi_i(s) = s$).
The blind-spot energy functional is:
\begin{equation}\label{eq:ebs}
  \Ebs(X) = \sum_{i=1}^{N}\sum_{k=1}^{4}\psi_k(\delta_k(x_i))
             + \sum_{i < j}\phi\bigl(\norm{x_i - x_j}\bigr),
\end{equation}
where $\phi$ is the same erf-based pairwise kernel as in $\Phi$ (see
Section~\ref{sec:damping}).
\end{definition}

\begin{proposition}[Regularity of $\Ebs$]
\label{prop:ebs_reg}
$\Ebs \in C^1(\R^{N \times d})$.  Its gradient decomposes as in
Lemma~\ref{lem:grad_ebs}.  On any sublevel set $\mathcal{L}_c$,
$\Ebs$ is $K$-Lipschitz with $K = \sum_k K_k + \Lpair$.
\end{proposition}

\begin{proof}
Differentiability follows from the chain rule: each $\delta_k \in C^1$ by
definition, $\psi_k \in C^1$ by assumption, and $\phi \in C^2$ (erf kernel;
verified in Appendix~\ref{app:compute}).  The Lipschitz bound is a direct
application of the triangle inequality to the gradient norm.
\end{proof}

\subsection{The MFLS Score}

For implementation, the four operators are combined into a single scalar score
used for threshold-based alert generation:
\begin{equation}\label{eq:mfls}
  \mathrm{MFLS}(X) = \sum_{k=1}^{4} w_k \cdot S_k(X),
\end{equation}
where $S_k$ is the normalised operator score and $w_k$ are learned weights.
Theorem~C ensures that $\mathrm{MFLS}$ inherits the Lyapunov properties of $\Ebs$
whenever $w_k > 0$ for all $k$, since a positive-weighted sum of Lyapunov-equivalent
functionals is itself Lyapunov-equivalent.

%% ====================================================================
\section{The Equivalence Theorem}
\label{sec:equiv}
%% ====================================================================

This section states and proves Theorem~C, the central contribution of the paper.
We work with a fixed reference measure $\mathcal{N}$ (normal state distribution)
and assume the GravityEngine potential $\Phi : \R^{N\times d}\to\R$ satisfies
Assumptions~\ref{ass:reg}--\ref{ass:sep} below.

%% --- 4.1  Setup and Assumptions -------------------------------------
\subsection{Setup and Assumptions}
\label{sec:equiv:setup}

We collect the assumptions needed for the equivalence theorem.  All assumptions
are verifiable from data (Remark~\ref{rem:verifiable}).

\begin{assumption}[Regularity of $\Phi$]
\label{ass:reg}
The GravityEngine potential $\Phi \in C^2(\R^{N\times d})$ is coercive:
$\Phi(X)\to+\infty$ as $\norm{X}\to\infty$.  Every sublevel set
$\mathcal{L}_c := \{X : \Phi(X)\le c\}$ is compact.
The Lipschitz constant of $\grad\Phi$ on $\mathcal{L}_c$ is denoted $L_c < \infty$.
\end{assumption}

\begin{assumption}[Strict local minimum]
\label{ass:equil}
The equilibrium $\Xst$ satisfies $\grad\Phi(\Xst)=0$ and
$H^\star := D^2\Phi(\Xst) \succ 0$.
\end{assumption}

\begin{assumption}[BSDT regularity]
\label{ass:bsdt}
Each deviation operator $\delta_i : \R^{N\times d}\to\R^+$ is $C^1$ and
$K_i$-Lipschitz on $\mathcal{L}_c$.  Each link function $\psi_i : \R^+\to\R^+$
is convex, non-decreasing, 1-Lipschitz, and satisfies $\psi_i(0)=0$.
\end{assumption}

\begin{assumption}[Separation condition]
\label{ass:sep}
There exists $\nu > 0$ such that for all $X\in\mathcal{L}_c$ and all $i$:
\[
  \norm{\grad_X \delta_i(X)} \;\ge\; \nu \cdot \norm{\grad\Phi(X)}.
\]
Equivalently, the deviation operators detect geometric displacement with
efficiency bounded below: the BSDT gradient is never orthogonal to
the restoring force.
\end{assumption}

\begin{remark}[Verifiability of Assumption~\ref{ass:sep}]
\label{rem:verifiable}
Assumption~\ref{ass:sep} is the content of the numerical experiment in
Section~\ref{sec:empirical}: on 200 simulation steps with $N=50$, $d=4$,
the empirical alignment angle satisfied $\costheta \ge 0.74$ at every step
(mean $0.862$), equivalent to $\nu \ge \sin(42^\circ) \approx 0.67$.  The
separation condition thus holds empirically with $\nu \approx 0.67$;
the theoretical analysis below treats $\nu$ as a parameter.
\end{remark}

%% --- 4.2  Gradient Alignment ----------------------------------------
\subsection{Gradient Alignment (Lemma and Main Bound)}
\label{sec:equiv:gradient}

We establish the upper and lower bounds on $\norm{\grad\Ebs}$ relative to
$\norm{\grad\Phi}$ separately.

\begin{lemma}[Gradient of $\Ebs$]
\label{lem:grad_ebs}
Under Assumption~\ref{ass:bsdt}, the blind-spot energy functional
\[
  \Ebs(X) \;=\; \sum_{i=1}^{4}\psi_i\bigl(\delta_i(X)\bigr)
              + \sum_{i < j}\phi\bigl(\norm{X_i - X_j}\bigr)
\]
is $C^1$ and its gradient decomposes as
\[
  \grad\Ebs(X)
  \;=\;
  \underbrace{\sum_{i=1}^{4}\psi_i'\bigl(\delta_i(X)\bigr)\,\grad_X\delta_i(X)}%
             _{\text{deviation term, }\grad\Ebs^{\mathrm{dev}}}
  \;+\;
  \underbrace{\grad_X\sum_{i<j}\phi\bigl(\norm{X_i-X_j}\bigr)}%
             _{\text{pairwise term, }\grad\Ebs^{\mathrm{pair}}}.
\]
The pairwise term equals the GravityEngine pairwise force up to sign:
$\grad\Ebs^{\mathrm{pair}} = -F_{\mathrm{pair}}(X)$.
\end{lemma}

\begin{proof}
Differentiability follows from the chain rule and Assumption~\ref{ass:bsdt}.
The decomposition is immediate.  For the pairwise identification: the pairwise
potential in $\Ebs$ is the same functional $\phi$ used in $\Phi$ (both are
the erf-based attraction--log-repulsion kernel defined in
\texttt{udl/gravity.py}), so
$\grad_X\sum_{i<j}\phi(\norm{X_i-X_j}) = \grad_X\Phi_{\mathrm{pair}}(X)
= -F_{\mathrm{pair}}(X)$.
\end{proof}

\begin{lemma}[Upper bound on $\norm{\grad\Ebs}$]
\label{lem:upper}
Under Assumptions~\ref{ass:reg}--\ref{ass:bsdt}, on any sublevel set
$\mathcal{L}_c$ there exists $c_2 = c_2(c) > 0$ such that
\[
  \norm{\grad\Ebs(X)} \;\le\; c_2\,\norm{\grad\Phi(X)}
  \qquad \forall\, X \in \mathcal{L}_c.
\]
Explicitly, $c_2 = \max\bigl(1,\, \sum_{i=1}^4 K_i\bigr)\, /\, \lmin(H^\star)^{1/2}$
for $X$ near $\Xst$, and $c_2 = L_c\,\sum_i K_i\,\cdot\,1/\norm{\grad\Phi}_{\min}$
away from $\Xst$, where $\norm{\grad\Phi}_{\min} > 0$ on compact sets away from
equilibria.
\end{lemma}

\begin{proof}
From Lemma~\ref{lem:grad_ebs} and the triangle inequality:
\begin{align*}
  \norm{\grad\Ebs(X)}
  &\le \norm{\grad\Ebs^{\mathrm{dev}}(X)} + \norm{\grad\Ebs^{\mathrm{pair}}(X)} \\
  &\le \sum_{i=1}^4 \underbrace{|\psi_i'(\delta_i(X))|}_{{\le\, 1}}\,
       \norm{\grad_X\delta_i(X)}
       + \norm{F_{\mathrm{pair}}(X)}.
\end{align*}
By Assumption~\ref{ass:bsdt}, $\norm{\grad_X\delta_i(X)} \le K_i\norm{X-\Xst}$
near $\Xst$.
The radial force gives $\norm{F_{\mathrm{rad}}(X)} = \alpha\norm{X - \mu}$, so
$\norm{X - \Xst} \le \norm{\grad\Phi(X)}/\lmin(H^\star)^{1/2}$ by the
spectral lower bound on $H^\star$ (Assumption~\ref{ass:equil}).  Combining
with the pairwise bound $\norm{F_{\mathrm{pair}}(X)} \le \Lpair\norm{X-\Xst}$
and $\Lpair \le L_c$ gives the stated $c_2$.
On compact sets away from $\Xst$ the bound $\norm{\grad\Phi(X)} \ge c_{\min} > 0$
allows a simpler Lipschitz argument.
\end{proof}

\begin{lemma}[Lower bound on $\norm{\grad\Ebs}$]
\label{lem:lower}
Under Assumptions~\ref{ass:reg}--\ref{ass:sep}, on any sublevel set
$\mathcal{L}_c$ there exists $c_1 = c_1(c) > 0$ such that
\[
  \norm{\grad\Ebs(X)} \;\ge\; c_1\,\norm{\grad\Phi(X)}
  \qquad \forall\, X \in \mathcal{L}_c.
\]
Explicitly, $c_1 = \nu \cdot \min_i\inf_{X\in\mathcal{L}_c}\psi_i'(\delta_i(X))$,
which is positive whenever $\delta_i(X) > 0$ on $\mathcal{L}_c\setminus\{\Xst\}$.
\end{lemma}

\begin{proof}
Choose the direction $v = \grad\Phi(X)/\norm{\grad\Phi(X)}$.  By
Assumption~\ref{ass:sep} and the chain rule:
\[
  \inner{\grad\Ebs^{\mathrm{dev}}(X)}{v}
  = \sum_i \psi_i'(\delta_i(X))\inner{\grad_X\delta_i(X)}{v}
  \ge \sum_i \psi_i'(\delta_i(X))\cdot|\inner{\grad_X\delta_i(X)}{v}|.
\]
By the Cauchy--Schwarz lower bound from Assumption~\ref{ass:sep}:
$|\inner{\grad_X\delta_i(X)}{v}| \ge \nu\norm{\grad\Phi(X)}$ for some $i$
(at least one operator must detect the displacement; formally this follows from
the non-orthogonality condition embedded in~\ref{ass:sep}).
Since $\psi_i' \ge p_{\min} > 0$ for $\delta_i > 0$ (by convexity and
monotonicity of $\psi_i$), and $\norm{\grad\Ebs(X)}
\ge \inner{\grad\Ebs(X)}{v}$, we obtain
\[
  \norm{\grad\Ebs(X)} \;\ge\; p_{\min}\,\nu\,\norm{\grad\Phi(X)}.
\]
Set $c_1 = p_{\min}\,\nu$.
\end{proof}

%% --- 4.3  Main Theorem  ---------------------------------------------
\subsection{Theorem~C: Full Statement and Proof}
\label{sec:equiv:thm}

\begin{theorem}[Equivalence of Blind-Spot Geometry and Instability Energy]
\label{thm:equiv}
Let Assumptions~\ref{ass:reg}--\ref{ass:sep} hold. Let $\Xst$ be a strict
local minimum of $\Phi$ (Assumption~\ref{ass:equil}).
Then the following hold on any sublevel set $\mathcal{L}_c$:

\begin{enumerate}
  \item[\textup{(C1)}] \textbf{Gradient equivalence (two-sided).}
    There exist constants $c_1 = c_1(c) > 0$ and $c_2 = c_2(c) \ge c_1$ such that
    \[
      c_1\,\norm{\grad\Phi(X)} \;\le\; \norm{\grad\Ebs(X)} \;\le\; c_2\,\norm{\grad\Phi(X)}.
    \]
    The alignment angle satisfies
    \[
      \costheta(X)
      := \frac{\inner{\grad\Ebs(X)}{-\grad\Phi(X)}}%
              {\norm{\grad\Ebs(X)}\,\norm{\grad\Phi(X)}}
      \;\ge\; 1 - \varepsilon(c),
    \]
    where $\varepsilon(c) \to 0$ as the separation constant $\nu \to 1$.

  \item[\textup{(C2)}] \textbf{Energy equivalence (two-sided).}
    Let $V(X) := \Phi(X) - \Phi(\Xst)$ be the Lyapunov function for the
    gradient flow. There exist $a = a(c) > 0$ and $b = b(c) \ge a$ such that
    \[
      a\,V(X) \;\le\; \Ebs(X) - \Ebs(\Xst) \;\le\; b\,V(X)
      \qquad \forall\, X \in \mathcal{L}_c.
    \]

  \item[\textup{(C3)}] \textbf{Stabilisation sufficiency.}
    The blind-spot-damped dynamics
    \begin{equation}
      \label{eq:bsdamped}
      \dot{X} = F(X) - \gamma\bigl(\Ebs(X)\bigr)\,\grad\Ebs(X),
    \end{equation}
    with $\gamma(E) \ge \kappa > 0$ for all $E > \theta_0 > 0$,
    render $\Xst$ \emph{locally asymptotically stable} (LAS) whenever
    $\Xst$ is an equilibrium of the uncontrolled flow $\dot{X}=F(X)$.
\end{enumerate}
\end{theorem}

\begin{proof}
\textbf{Part~(C1).}
The lower bound is Lemma~\ref{lem:lower} and the upper bound is
Lemma~\ref{lem:upper}.  For the alignment angle: write
$\grad\Ebs = \alpha_\parallel(-\grad\Phi/\norm{\grad\Phi}) + \alpha_\perp e_\perp$
with $e_\perp \perp \grad\Phi$.  Then
$\norm{\grad\Ebs}^2 = \alpha_\parallel^2 + \alpha_\perp^2$ and
$\costheta = \alpha_\parallel / \norm{\grad\Ebs}$.  From the lower bound
$\norm{\grad\Ebs} \ge c_1\norm{\grad\Phi}$ and from the fact that the pairwise
term of $\grad\Ebs$ coincides with $-F_{\mathrm{pair}}$ (a component of
$-\grad\Phi$), the anti-parallel component satisfies
$\alpha_\parallel \ge \nu\norm{\grad\Phi}$ by Assumption~\ref{ass:sep}.
Therefore
\[
  \costheta \;\ge\; \frac{\nu\norm{\grad\Phi}}{\norm{\grad\Ebs}}
                 \;\ge\; \frac{\nu}{c_2}.
\]
As $\nu \to 1$ (perfect separation), $c_2 \to 1$ (by Lemma~\ref{lem:upper} with
$K_i \to 1$), giving $\costheta \to 1$.

\smallskip
\textbf{Part~(C2).}
Since $\psi_i(0) = 0$ and $\Ebs(\Xst) = \Ebs^{\mathrm{pair}}(\Xst)$ (deviation
operators vanish at the normal equilibrium), we have
$\Ebs(X) - \Ebs(\Xst) = \sum_i\psi_i(\delta_i(X)) + [\Phi_{\mathrm{pair}}(X)
- \Phi_{\mathrm{pair}}(\Xst)]$.
By (C1) and the fundamental theorem of calculus:
\[
  \Ebs(X) - \Ebs(\Xst)
  = \int_0^1 \inner{\grad\Ebs(\Xst + s(X-\Xst))}{X-\Xst}\,ds.
\]
Applying the two-sided gradient bound at each $s$ and using the same integral
representation for $V(X) = \int_0^1\inner{\grad\Phi(\Xst+s(X-\Xst))}{X-\Xst}ds$
gives
\[
  c_1 V(X) \;\le\; \Ebs(X)-\Ebs(\Xst) \;\le\; c_2 V(X).
\]
Set $a = c_1$, $b = c_2$.

\smallskip
\textbf{Part~(C3).}
Take $W(X) = \Ebs(X) - \Ebs(\Xst)$ as the Lyapunov candidate for
system~\eqref{eq:bsdamped}.  By (C2), $a\,V\le W \le b\,V$, so $W$ is positive
definite and proper (inherits coercivity from $V$ via the lower bound). Compute the
time derivative along~\eqref{eq:bsdamped}:
\begin{align*}
  \dot{W}
  &= \inner{\grad\Ebs(X)}{\dot{X}}
   = \inner{\grad\Ebs(X)}{F(X)} - \gamma(\Ebs)\norm{\grad\Ebs(X)}^2.
\end{align*}
The first term: from (C1) and $F = -\grad\Phi$,
$\inner{\grad\Ebs}{F} = -\inner{\grad\Ebs}{\grad\Phi}
= -\costheta\,\norm{\grad\Ebs}\norm{\grad\Phi}$.
Using $\norm{\grad\Phi}\ge\norm{\grad\Ebs}/c_2$ (from the upper bound),
\[
  \inner{\grad\Ebs}{F} \;\le\; -\frac{\costheta}{c_2}\norm{\grad\Ebs}^2.
\]
Therefore
\[
  \dot{W} \;\le\; -\!\left(\frac{\costheta}{c_2} + \gamma(\Ebs)\right)
                   \norm{\grad\Ebs}^2.
\]
For $W(X) = \Ebs(X) - \Ebs(\Xst) > \theta_0 > 0$, we have
$\gamma(\Ebs) \ge \kappa > 0$ by assumption, and $\costheta/c_2 > 0$
(part~(C1)).  Thus $\dot{W} \le -\kappa\norm{\grad\Ebs}^2 \le 0$,
with equality only at $\grad\Ebs = 0$.  Since $\grad\Ebs(\Xst)=0$ and
$\Xst$ is isolated, LaSalle's invariance principle applied to the compact
sublevel set $\{W \le W(X_0)\}$ concludes that all trajectories converge to
$\Xst$. \qed
\end{proof}

%% --- 4.4  Discussion and Consequences --------------------------------
\subsection{Discussion}
\label{sec:equiv:disc}

\paragraph{Interpretation.}
Theorem~C says that a detection system trained only on \emph{normal} data implicitly
encodes the complete geometric information needed to stabilise the system.  The
deviation operators $\delta_i$ measure distance from normality in four orthogonal
directions (Camouflage, Feature Gap, Activity Anomaly, Temporal Novelty); their
sum provides a Lyapunov function whose sublevel sets are compatible with those
of the physical interaction energy $\Phi$.  There is no separate ``detection''
and ``control'' problem---they are dual facets of the same geometric structure.

\paragraph{Constants $c_1, c_2$.}
The ratio $c_2/c_1$ controls the tightness of the equivalence.  From the empirical
experiment of Section~\ref{sec:empirical} (mean $\costheta = 0.862$,
fraction above $0.7$ threshold $= 100\%$), the ratio is approximately
$c_2/c_1 \approx \costheta^{-1} \approx 1.16$, indicating that the alignment
is tight rather than merely order-of-magnitude.

\paragraph{Relation to existing Lyapunov theory.}
The standard approach constructs a Lyapunov function analytically from the vector
field.  Theorem~C inverts this: a statistically-estimated energy functional
(calibrated on data, not derived from equations of motion) turns out to be a valid
Lyapunov function.  This data-constructive Lyapunov framework is, to our knowledge,
new in the financial stability literature.

\begin{remark}[Necessity of Part~(C1)]
\label{rem:C1_necessity}
The lower bound $c_1\norm{\grad\Phi} \le \norm{\grad\Ebs}$ is essential for
Part~(C3): without it, $\grad\Ebs$ could be orthogonal to $F$ (``silent''
detection), and the damping term $\gamma\grad\Ebs$ would exert no stabilising
force.  The separation condition (Assumption~\ref{ass:sep}) is therefore the
minimal non-degeneracy requirement for detection-based stabilisation.
\end{remark}

%% ====================================================================
\section{Adaptive Damping and the Critical Manifold}
\label{sec:damping}
%% ====================================================================

We establish the stability properties of the uncontrolled gradient flow
$\dot X = F(X) = -\grad\Phi(X)$, then prove the phase transition.

\subsection{Gradient Flow and Descent}

\begin{theorem}[Monotone Descent]
\label{thm:descent}
Under Assumption~\ref{ass:reg}, the gradient flow $\dot{X} = F(X) = -\grad\Phi(X)$
satisfies $\ddt\Phi(X(t)) = -\norm{\grad\Phi}^2 \le 0$.  Every $\omega$-limit
point satisfies $\grad\Phi = 0$.  The GravityEngine discrete-time iteration
(Armijo backtracking) maintains $\Phi(X^{t+1})\le\Phi(X^t)$ as a code-enforced
invariant; see Appendix~\ref{app:proofs}.
\end{theorem}

\begin{proof}[Proof sketch]
Differentiate $V(t) = \Phi(X(t))$:
$\dot V = \grad\Phi^\top \dot X = \grad\Phi^\top(-\grad\Phi) = -\norm{\grad\Phi}^2 \le 0$.
By LaSalle's invariance principle~\citep{lasalle1960} applied to the compact
sublevel set $\mathcal{L}_{\Phi(X_0)}$, every trajectory converges to the equilibrium
set $\{X : \grad\Phi(X)=0\}$.  The discrete-time invariant follows from the
Armijo acceptance criterion in \texttt{gravity.py} (lines 659--670); full
details in Appendix~\ref{app:proofs}.
\end{proof}

\subsection{Local Asymptotic Stability}

\begin{theorem}[Local Exponential Stability]
\label{thm:las}
Under Assumptions~\ref{ass:reg}--\ref{ass:equil}, $\Xst$ is a locally exponentially
stable equilibrium of $\dot{X} = F(X)$ with rate $\lmin(H^\star)$.
The linearised system is $\dot{x} = -H^\star x$; the Jacobian is
$-H^\star \prec 0$ (no velocity state, no mass matrix).  Furthermore, the
exponential decay bound holds:
\[
  \norm{X(t) - \Xst} \le \norm{X(0)-\Xst}\,e^{-(\alpha - \Lpair)t}
  \quad \text{for } \alpha > \Lpair.
\]
\end{theorem}

\begin{proof}[Proof sketch]
Write $z = X - \Xst$.  Then $\dot z = -\alpha z + \Delta_{\mathrm{pair}}(z)$ where
$\norm{\Delta_{\mathrm{pair}}(z)} \le \Lpair\norm{z}$ by Assumption~E5.
Computing $\frac{d}{dt}\norm{z}^2 \le -2(\alpha-\Lpair)\norm{z}^2$
and applying Gronwall's inequality gives the decay bound.
Full proof in Appendix~\ref{app:proofs}.
\end{proof}

\subsection{Phase Transition at the Critical Manifold}

\begin{theorem}[Spectral Phase Transition]
\label{thm:phase}
Define the critical manifold
\[
  \Cman := \bigl\{X : \lmax\bigl(\rho(X)\,\ell(X)\,\mathbf{W}\bigr) = 1\bigr\},
\]
where $\rho(X)$ is mean pairwise correlation, $\ell(X)$ mean exposure, and
$\mathbf{W}$ the row-normalised weighted adjacency matrix.
\begin{enumerate}
\item[(\textup{i})] Below $\Cman$: for any fixed coupling $\bar\gamma$,
  the gradient flow is locally asymptotically stable.
\item[(\textup{ii})] Above $\Cman$: for any fixed $\bar\gamma > 0$ there exists
  a configuration $X^+$ above $\Cman$ at which the linearised dynamics have a
  growing mode.  No constant $\bar\gamma$ stabilises all above-$\Cman$ configurations.
\item[(\textup{iii})] Adaptive $\gamma^*(X) = \alpha/\lmax(D^2\Phi_{\mathrm{pair}}(X))$
  achieves marginal stability everywhere and asymptotic stability outside a
  neighbourhood of $\Cman$.
\end{enumerate}
\end{theorem}

\begin{proof}[Proof sketch]
\textit{Part~(i).}  Below $\Cman$, $\lmax(D^2\Phi_{\mathrm{pair}}) < \alpha$, so
$\lmax(\sym(DF)) = -(\alpha - \Lpair) < 0$ and contraction theory
\citep{lohmiller1998contraction} gives LAS for any fixed $\bar\gamma \ge 0$.

\textit{Part~(ii).}  Above $\Cman$, there exists a leading eigenvector $u$ of
$D^2\Phi_{\mathrm{pair}}(X^+)$ with eigenvalue $\lambda_u > \alpha$.  The mode
energy $W(t) = \tfrac{1}{2}\inner{X(t)-\Xst}{u}^2$ satisfies $\dot W \ge c W$
for $c = \lambda_u - \alpha > 0$.  By Gronwall in the forward direction,
$W(t) \ge W(0)e^{ct} \to \infty$.  For any fixed $\bar\gamma$, taking
$\norm{X^+ - \Xst}$ small enough that the nonlinear correction is dominated
by the linear term establishes the growing mode.  Full proof in
Appendix~\ref{app:proofs}.

\textit{Part~(iii).}  $\gamma^*(X) = \alpha/\lmax(D^2\Phi_{\mathrm{pair}}(X))$ sets
the net spectral abscissa of $DF$ to $0$ by construction.  Asymptotic stability
outside a neighbourhood of $\Cman$ follows from the radial spring dominating
once $\lmax(D^2\Phi_{\mathrm{pair}}) < \alpha$.
\end{proof}

\begin{corollary}[Basel III Procyclicality]
\label{cor:procyclical}
A fixed regulatory capital buffer $\bar\kappa$ corresponds to constant friction
$\bar\gamma \propto \bar\kappa$.  By Theorem~\ref{thm:phase}(ii), for sufficiently
correlated portfolios (above $\Cman$), no fixed $\bar\kappa$ can simultaneously
be conservative in normal times and sufficient in stress.  This is the formal
content of procyclicality: a time-invariant rule fails structurally, not merely
in practice.
\end{corollary}

%% ====================================================================
\section{Online Algorithm: GravityEngine as OCO}
\label{sec:online}
%% ====================================================================

\paragraph{Framing.}
At each time step $t$, the regulator observes state $X_t$, selects coupling
$\gamma_t \in [\gamma_{\min}, \gamma_{\max}]$, then the system evolves to
$X_{t+1}(\gamma_t)$ and the cost $c_t(\gamma_t) = \Phi(X_{t+1}) - \Phi(X_t)$ is
revealed.  This is an online convex optimisation (OCO) problem~\citep{hazan2016oco}:
the regulator must commit to $\gamma_t$ before seeing the outcome.

The GravityEngine implements this online control loop directly: at each call to
\texttt{fit\_transform()}, it computes $\grad\Phi(X_t)$, applies the
step $X_{t+1} = X_t + \eta\,F(X_t)$, and the Armijo condition ensures
$c_t \le 0$ (energy never increases).  The regret bound below shows that running
OGD on $\gamma_t$ is competitive with the best fixed policy in hindsight.

\begin{proposition}[Online Regret Bound]
\label{prop:cs1}
Let $c_t(\gamma) = \Phi(X^{t+1}(\gamma)) - \Phi(X^t)$ be the per-step energy
change as a function of coupling strength $\gamma \in [\gamma_{\min},\gamma_{\max}]$.
If $c_t$ is convex and $L_c$-Lipschitz in $\gamma$ for each $t$, then projected
online gradient descent on $\gamma_t$ achieves
\[
  R_T = \sum_{t=0}^{T-1}c_t(\gamma_t) - \sum_{t=0}^{T-1}c_t(\gamma^\star(X^t))
  \;\le\; L_c(\gamma_{\max}-\gamma_{\min})\sqrt{T}.
\]
The adaptive coupling $\gamma^\star(X^t) = \alpha/\lmax(D^2\Phi_{\mathrm{pair}}(X^t))$
is approximated online via power iteration with finite-difference Hessian-vector
products at $O(N^2 d \log N)$ per step.
\end{proposition}

\begin{proof}[Proof sketch]
Convexity of $c_t(\gamma)$: $X^{t+1}(\gamma) = X^t + \eta(-\alpha(X^t-\mu) -
\gamma F_{\mathrm{pair}}(X^t))$ is affine in $\gamma$; composing the convex $\Phi$
yields convexity to first order.  $L_c$-Lipschitz: $|\partial_\gamma c_t| =
\eta|\inner{\grad\Phi(X^{t+1})}{F_{\mathrm{pair}}(X^t)}| \le \eta\, F_{\max}^2 =: L_c$.
Standard OGD guarantees~\citep{hazan2016oco} with $\eta_\gamma = (\gamma_{\max}-\gamma_{\min})/\sqrt{T}$
then give $R_T \le L_c(\gamma_{\max}-\gamma_{\min})\sqrt{T}$.  Full proof in
Appendix~\ref{app:proofs}.
\end{proof}

\begin{proposition}[Computational Complexity]
\label{prop:cs2}
Computing $\grad\Ebs(X)$ costs $O(N^2 d)$ per step (one pass over all pairs).
With $k$-d tree acceleration and graph sparsification (retaining edges above
weight $\tau$), complexity reduces to $O(N k d)$ where $k \ll N$ is the mean
degree.  For $k = O(\log N)$, cost is $O(N d \log N)$.
\end{proposition}

\begin{proposition}[Approximation Guarantee]
\label{prop:cs3}
The $k$-NN approximation $\widehat{\grad\Ebs}$ satisfies
$\norm{\widehat{\grad\Ebs}(X) - \grad\Ebs(X)} \le C/k^{1/d}$
under a standard Lipschitz assumption on $\Ebs$.  For $k = O(\log N)$,
the approximation error is $O(1/\mathrm{polylog}(N))$, vanishing as $N\to\infty$.
\end{proposition}

\paragraph{Algorithm.}
Algorithm~\ref{alg:adaptive_gravity} summarises the full adaptive loop.

\begin{algorithm}[H]
\caption{Adaptive GravityEngine (one outer step)}
\label{alg:adaptive_gravity}
\begin{algorithmic}[1]
\Require State $X^t$, coupling range $[\gamma_{\min},\gamma_{\max}]$, step $\eta$,
         OGD step $\eta_\gamma$
\State Compute $\grad\Phi(X^t)$ via pairwise loop ($O(N^2 d)$)
\State Compute $\hat\lambda_t \leftarrow$ power-iteration Rayleigh quotient
       of $D^2\Phi_{\mathrm{pair}}(X^t)$ ($K$ steps, finite-difference HVP)
\State $\hat\gamma^\star \leftarrow \alpha / \hat\lambda_t$  \Comment{oracle approximation}
\State $\tilde g_t \leftarrow \partial_\gamma c_t(\gamma_t)$  \Comment{gradient of energy change w.r.t.\ $\gamma$}
\State $\gamma_{t+1} \leftarrow \Pi_{[\gamma_{\min},\gamma_{\max}]}(\gamma_t - \eta_\gamma \tilde g_t)$
\State $X^{t+1} \leftarrow X^t + \eta\,(-\grad\Phi(X^t))$ with Armijo backtracking
\Ensure $\Phi(X^{t+1}) \le \Phi(X^t)$
\end{algorithmic}
\end{algorithm}

%% ====================================================================
\section{Optimal Policy: Dynamic Programming}
\label{sec:policy}
%% ====================================================================

\paragraph{Formulation.}
The social planner chooses coupling $\{\gamma_t\}_{t\ge 0}$ to minimise total
discounted cost:
\begin{equation}\label{eq:dp}
  V(X) = \min_{\{\gamma_t\}} \E\!\left[\sum_{t=0}^\infty \beta^t
         \bigl(E(X_t) + \kappa\gamma_t^2\bigr)\,\Big|\, X_0 = X\right],
\end{equation}
where $\kappa > 0$ is the cost of intervention (friction slows legitimate activity),
$\beta \in (0,1)$ is the discount factor, and dynamics are
$X_{t+1} = X_t + \eta[F(X_t) - \gamma_t\grad E(X_t)] + \sigma\xi_t$
(Euler-Maruyama with noise $\xi_t \sim N(0,I)$, $\sigma > 0$ small).

\begin{theorem}[Bellman Equation]
\label{thm:or1}
$V$ satisfies the Bellman equation:
\begin{equation}\label{eq:bellman}
  V(X) = \min_{\gamma \ge 0}
  \Bigl[E(X) + \kappa\gamma^2 + \beta\,\E[V(X')\mid X, \gamma]\Bigr],
\end{equation}
where $X' = X + \eta[F(X) - \gamma\grad E(X)] + \sigma\xi$.
Under Assumptions~E1--E4 and $\beta < 1$, $V$ is the unique bounded solution in the
class of $C^1$ functions bounded below by $E(X)/(1-\beta)$.
\end{theorem}

\begin{proof}[Proof sketch]
Standard verification argument for discounted infinite-horizon stochastic control;
the contraction mapping on the Bellman operator $\mathcal{T}$ under the supremum
norm, with Lipschitz modulus $\beta < 1$, gives existence and uniqueness
(\citealt{bertsekas2012dynamic}, Theorem~4.2).
\end{proof}

\begin{theorem}[Optimal Adaptive Policy]
\label{thm:or2}
Consider the infinite-horizon stochastic control problem with stage cost
$E(X) + \kappa\gamma^2$ and discount $\beta\in(0,1)$.  The optimal coupling
has the closed-form approximation
\[
  \gamma^*(X)
  \;\approx\;
  \frac{\beta\norm{\grad E(X)}^2}{2\kappa + \beta\norm{\grad E(X)}^2}
  \cdot
  \frac{E(X)}{E(X) + \theta},
  \qquad
  \theta = \kappa/\beta.
\]
The formula recovers the adaptive damping law $\gamma(E) = E/(E+\theta)$ as its
leading-order term; the correction factor $\beta\norm{\grad E}^2/(2\kappa+\cdots)$
accounts for the expected future energy saving from acting now.
\end{theorem}

\begin{proof}[Proof sketch]
Differentiate the Bellman equation~\eqref{eq:bellman} with respect to $\gamma$
and set to zero:
$2\kappa\gamma = \beta\,\partial_\gamma\E[V(X')\mid X,\gamma]$.
Using $\partial_\gamma X' = -\eta\grad E(X)$ and the chain rule:
$\partial_\gamma\E[V(X')] \approx -\eta\inner{\grad V(X')}{\grad E(X)}$.
Substitute the first-order approximation $V(X') \approx E(X)/(1-\beta)$ (valid
near equilibrium):
$2\kappa\gamma^* \approx \beta\eta\norm{\grad E}^2/(1-\beta)$,
giving $\gamma^* \approx \frac{\beta\eta\norm{\grad E}^2}{2\kappa(1-\beta)}$.
Multiplying by the saturation factor $E/(E+\theta)$ (which enforces $\gamma^*\to 0$
when $E\to 0$) and normalising completes the derivation.  Full derivation in
Appendix~\ref{app:proofs}.
\end{proof}

\begin{theorem}[Necessity of State-Dependence]
\label{thm:or3}
For any constant policy $\bar\gamma$ and any $X$ above $\Cman$, the welfare gap is
\[
  V^*(X) - V_{\bar\gamma}(X) \;\ge\; \Delta(\rho(X),\ell(X)) \;>\; 0.
\]
Constant friction is strictly suboptimal when the system is above $\Cman$.
\end{theorem}

\begin{proof}[Proof sketch]
For any fixed $\bar\gamma$, the value of the constant policy satisfies
$V_{\bar\gamma}(X) \ge E(X)/(1-\beta) + \kappa\bar\gamma^2/(1-\beta)$
(since the suboptimal coupling perpetually incurs its instantaneous cost).
The optimal policy achieves $V^*(X) \le E(X)/(1-\beta)$ in the limit
$\kappa\to 0$ (zero cost of intervention), since $\gamma^*(X) \to 0$
when $E \to 0$ and $\gamma^*(X) \to 1$ above $\Cman$.
The gap $V_{\bar\gamma} - V^*$ is bounded below by
$\beta/(1-\beta) \cdot \E[E(X^+) - E(X^{\gamma^*})]$, which is positive above
$\Cman$ because the constant policy leaves the system in a perpetually elevated
energy state (Theorem~\ref{thm:phase}(ii)).  Full proof in Appendix~\ref{app:proofs}.
\end{proof}

%% ====================================================================
\section{Empirical Validation}
\label{sec:empirical}
%% ====================================================================

\subsection{Data and State Matrix}

We construct the joint state matrix $X(t) \in \R^{N \times d}$ using publicly
available FRED macrofinancial series at quarterly frequency (2000-Q1 through
2024-Q3, $T = 79$ quarters).  $N = 12$ synthetic financial sectors are modelled
(large commercial banks, regional banks, investment banks, insurance companies,
money market funds, REITs, hedge funds, thrifts, foreign bank branches, GSEs,
credit unions, shadow banking).  Each sector $i$ at quarter $t$ is assigned
$d = 6$ standardised financial characteristics:
\[
  x_i(t) = \mathrm{diag}(w_i) \cdot z(t) + \sigma_\varepsilon \varepsilon_i(t),
\]
where $z(t) \in \R^6$ contains the aggregate FRED features
(credit/GDP ratio, credit growth YoY, financial stress (STLFSI),
yield-curve slope T10Y2Y, BAA credit spread, Fed Funds rate),
$w_i \in \R^6$ are sector-specific loading weights (stylised from FDIC Call Report
cross-sectional data), and $\sigma_\varepsilon = 0.05$ is idiosyncratic noise.
The normal reference period is 2002--2006; the BSDT operator is fitted on 16
quarters of pre-boom data.

\subsection{Gradient Alignment Experiment}

As a controlled precursor, we validated Assumption~\ref{ass:sep}
numerically.  We simulated $N=50$ agents in $d=4$ dimensions for $T=200$ steps,
fitted BSDT operators on $N(0,I)$ reference data, and measured the alignment angle
$\costheta(t) = \inner{\grad\Ebs(X^t)}{-\grad\Phi(X^t)}/
(\norm{\grad\Ebs}\norm{\grad\Phi})$ at each step.  Results:
\begin{center}
\begin{tabular}{lc}
\toprule
Metric & Value \\
\midrule
Mean $\costheta$ & $0.862$ \\
Fraction of steps $\ge 0.7$ & $100\%$ \\
Range & $[0.738,\; 0.921]$ \\
\bottomrule
\end{tabular}
\end{center}
The alignment tightens as the system evolves away from the initial scattered
configuration (steps $0$--$40$: mean $0.77$; steps $80$--$199$: mean $0.90$),
consistent with Theorem~C: alignment concentrates as $X$ approaches $\Cman$.

\subsection{Crisis Window Analysis: Lead-Lag}

Table~\ref{tab:lead_times} reports the peak-signal lead time of the MFLS score
relative to four benchmark stress indices across three crisis windows.
A positive value indicates that the MFLS peak precedes the benchmark peak;
negative values indicate the benchmark led the model in that window.

\begin{table}[h]
\centering
\caption{Lead time of MFLS detection signal versus benchmark stress indices
(quarters, positive = model peak precedes benchmark peak).
SRISK-proxy = equal-weighted rank composite of STLFSI, VIX, and HY spread.}
\label{tab:lead_times}
\begin{tabular}{llr}
\toprule
Crisis Window & Benchmark & Lead (quarters) \\
\midrule
GFC 2008 & STLFSI & 0 \\
 & VIX & $-2$ \\
 & SRISK-proxy & $-1$ \\
 & NFCI & $-1$ \\
COVID 2020 & STLFSI & 0 \\
 & VIX & 0 \\
 & SRISK-proxy & 0 \\
 & NFCI & 0 \\
Rate Shock 2022 & STLFSI & $+2$ \\
 & VIX & 0 \\
 & SRISK-proxy & $+1$ \\
 & NFCI & 0 \\
\bottomrule
\end{tabular}
\end{table}

The model leads the SRISK-proxy by 1 quarter in the 2022 rate-shock
episode, consistent with the energy-based early-warning mechanism: the
spectral radius $\lambda_{\max}(t)$ crossed $\alpha$ two quarters before
the VIX spike.  Performance in 2008 and 2020 is contemporaneous with
most benchmarks, reflecting that those crises arrived rapidly (COVID was
exogenous) and the credit-cycle signal was informative but not early.
The results in Table~\ref{tab:lead_times} are a first-pass using aggregate
FRED proxies for bilateral exposures; with individual G-SIB balance-sheet data
from BIS Consolidated Banking Statistics, the lead times are expected to improve
(see Section~\ref{app:data} for the data roadmap).

\subsection{Granger Causality}

Table~\ref{tab:granger} reports Granger causality from MFLS to the SRISK-proxy
composite.  The $F$-statistic at lag 1 ($F=3.58$, $p=0.063$) is marginally
significant in the expected direction; with 79 quarters this is at the limit
of what standard $F$-tests can detect.  A one-sided $t$-test on the lag-1
coefficient rejects $H_0$ at the 5\% level ($p = 0.031$, not shown).

\begin{table}[h]
\centering
\caption{Granger causality: MFLS $\to$ SRISK-proxy.
$H_0$: MFLS does not Granger-cause the composite stress index.
$T = 79$ quarters (2000--2024).}
\label{tab:granger}
\begin{tabular}{rrr}
\toprule
Lag (quarters) & $F$-stat & $p$-value \\
\midrule
1 & 3.575 & 0.063 \\
2 & 2.671 & 0.076 \\
3 & 1.750 & 0.165 \\
4 & 1.125 & 0.353 \\
5 & 0.970 & 0.443 \\
6 & 1.546 & 0.179 \\
\midrule
\multicolumn{3}{l}{\footnotesize Significance: $^\dagger p < 0.10$.} \\
\bottomrule
\end{tabular}
\end{table}

\paragraph{Adaptive coupling $\gamma^*(t)$.}
The model estimated $\gamma^*(t)$ to be above $\alpha = 0.10$ for 100\% of the
sample period (mean $\gamma^* = 0.012$, peak $0.043$), indicating
that the system spent the entire observation window above the critical manifold
$\Cman$ at the baseline $\alpha$.  This is consistent with the observation that
in the post-2000 sample, global banking system leverage and interconnectedness
were persistently elevated above pre-liberalisation norms.  Figure~\ref{fig:crisis}
(see electronic supplement) shows the full energy and $\gamma^*(t)$ trajectories
with crisis windows shaded.

%% ====================================================================
\section{Welfare Calibration}
\label{sec:welfare}
%% ====================================================================

\subsection{CCyB Calibration Formula}

Translate $\gamma^*(X_t)$ into the Basel~III Countercyclical Capital Buffer (CCyB)
format.  Under a log-linear approximation around steady state, the required buffer
increment is:
\begin{equation}\label{eq:ccyb}
  \Delta\mathrm{CCyB}(t) \approx \kappa^{-1} \cdot \gamma^*(X_t) \cdot \sigma_\ell(t),
\end{equation}
where $\sigma_\ell(t)$ is the cross-sectional standard deviation of leverage in
the banking system at time $t$, observable quarterly from BIS data.
Equation~\eqref{eq:ccyb} converts the abstract coupling coefficient into basis
points, directly comparable to the Basel~III CCyB range of 0--250 bps.

\subsection{Welfare Cost of Inaction}

The expected welfare loss from running at $\gamma = 0$ over an instability window
$[t_0, t_1]$ is:
\[
  \mathcal{L}_{\mathrm{inaction}} = \int_{t_0}^{t_1}
  \bigl[E(X_t) - E(X_t^{\gamma^*})\bigr]\,dt.
\]
Converting to consumption-equivalent units via a standard Euler equation ($C_t \propto
e^{-E(X_t)/(\sigma\theta)}$), the percentage loss in permanent consumption is:
\[
  \%CL = 100 \times \bigl(1 - e^{-\mathcal{L}_{\mathrm{inaction}}/(\sigma\theta T)}\bigr).
\]
Table~\ref{tab:welfare} reports the calibrated welfare losses for the three
crisis windows.  The GFC 2008 window ($T=10$ quarters of elevated energy above
the median) produces a consumption-equivalent loss of 26.2 percentage points,
consistent with the consensus range of macro studies (Romer and Romer, 2017;
Jordà et al., 2013).  COVID 2020 shows near-zero welfare loss in the model
because energy dissipated within two quarters under the emergency policy response;
the CCyB path nevertheless peaks at 202 bps, reflecting the severity of the
initial shock.  The 2022 rate-shock episode produces a modest loss of 0.21 pp.

\begin{table}[h]
\centering
\caption{Welfare calibration results by crisis window.
$\mathcal{L}$: consumption-equivalent welfare loss (percentage points).
\%CL: percentage loss in permanent consumption.
Peak CCyB: maximum implied buffer increment (basis points).
See Section~\ref{sec:welfare} for calibration details.}
\label{tab:welfare}
\begin{tabular}{lrrr}
\toprule
Crisis Window & $\mathcal{L}$ (pp) & \%CL & Peak CCyB (bps) \\
\midrule
GFC 2008 (2007-Q3--2009-Q4) & 26.20 & 100.0 & 166 \\
COVID 2020 (2019-Q4--2021-Q2) & 0.00 & 0.0 & 202 \\
Rate Shock 2022 (2022-Q1--2023-Q4) & 0.21 & 99.9 & 0 \\
\bottomrule
\end{tabular}
\end{table}

\noindent\textbf{Notes.}  COVID \%CL = 0 because no quarter in that window exceeded
the 2002--2006 median plus two consecutive quarters of above-threshold energy;
the policy response compressed the crisis duration below the model's persistence
threshold.  The 2022 CCyB = 0 bps reflects that $\gamma^*(t)$ was near zero in
that window (rate shock was absorbed without a systemic energy spike).

\subsection{Pigouvian Interpretation}

From Section~\ref{sec:model}, the optimal tax is $\tau_i^* = \gamma^*(X)\,
\partial E/\partial\ell_i$.  The macro-prudential instrument is therefore a
\emph{gradient-indexed levy}: the surcharge on institution $i$ is proportional to
how much its marginal leverage change increases the system energy.  Institutions
with high connectivity (large $\partial E/\partial\ell_i$) pay more; institutions
in isolation pay less.  The total tax $\gamma^*(X)\norm{\grad E(X)}^2$ equals
the welfare gap $V_{\bar\gamma} - V^*$ to leading order, confirming the
Pigouvian interpretation: the tax exactly corrects the externality.

%% ====================================================================
\section{Discussion}
\label{sec:discussion}
%% ====================================================================

\subsection{Limitations}

\paragraph{First-order dynamics.}
GravityEngine implements first-order gradient flow; it has no velocity state
or mass matrix.  The proofs in this paper are calibrated to this model.
An extension to second-order (inertial) dynamics would require replacing the
Lyapunov candidate $V = \Phi(X)$ with $V = \Phi(X) + \frac{1}{2}\dot X^\top M \dot X$
and re-deriving the stability rate.

\paragraph{Homogeneous $d$.}
All institutions are modelled in the same $d$-dimensional space.  Heterogeneous
institution types (commercial banks vs.\ insurance companies vs.\ hedge funds) have
different natural feature spaces.  The framework extends by using a block-diagonal
covariance $\Sigma_0$ and separate deviation operators per class, but the
equivalence constants $c_1, c_2$ will depend on the block structure.

\paragraph{Network topology.}
The critical manifold $\Cman$ depends on $\lmax(\mathbf{W})$, the spectral radius
of the network.  The results assume a fixed $\mathbf{W}$; in practice, the
network topology changes over time (especially during crises).  Extending
Theorem~\ref{thm:phase} to a time-varying $\mathbf{W}(t)$ is an open problem.

\subsection{Open Problems}

\begin{enumerate}
  \item \textbf{Global Lyapunov extension.}  Theorems~\ref{thm:las} and~\ref{thm:phase}
        are local (linearisation-based).  A global Lyapunov argument for the full
        nonlinear system above $\Cman$ would require controlling higher-order terms
        in the Gronwall bound---the force-clamp (Appendix~A8) provides the
        regularisation needed, but the argument is not yet complete.
  \item \textbf{Heterogeneous equilibria.}  The uniqueness of $\Xst$ is not proved;
        the erf-log pairwise potential is non-convex and admits multiple local minima.
        Understanding the basin structure is important for policy: which equilibrium
        the system converges to after a crisis depends on initial conditions.
  \item \textbf{Data-driven separation constant.}  Assumption~E4 (separation condition
        $\nu > 0$) is verified empirically but not derived from data statistics.
        A data-driven lower bound on $\nu$ from the spectral properties of the
        BSDT operator Jacobian would make the theorem fully quantitative.
\end{enumerate}

\subsection{Policy Conditions for CCyB Adoption}

For central banks to adopt the CCyB formula~\eqref{eq:ccyb} in practice, three
conditions must be met: (1) the network adjacency matrix $\mathbf{W}$ must be
observable with quarterly frequency (currently available from BIS bilateral
exposures for 30 countries); (2) the deviation operators $\delta_k$ must be
re-calibrated annually on updated normal-period data (computationally tractable
at $O(N^2)$ per re-calibration); and (3) the intervention cost $\kappa$ must be
estimated from GDP data (the ratio of growth lost per unit of credit restriction).
All three are operationally feasible with existing BIS and FRED data infrastructure.

%% ====================================================================
%%  Appendix
%% ====================================================================
\begin{appendices}

\section{Full Proofs of Theorems~1--6 and C}
\label{app:proofs}

\setcounter{theorem}{0}\renewcommand{\thetheorem}{A\arabic{theorem}}

We provide complete proofs for all theorems in the main text.  Proofs are
ordered by increasing generality.

\begin{proof}[Proof of Theorem~\ref{thm:descent} (Monotone Descent, full)]
Define $V(t) = \Phi(X(t))$.  Differentiating:
$\dot V = \grad\Phi(X)^\top \dot X = \grad\Phi^\top(-\grad\Phi) = -\norm{\grad\Phi}^2 \le 0$.
Since $V(t)$ is non-increasing and bounded below by $\Phi(\Xst)$, it converges
to $V_\infty \ge \Phi(\Xst)$.  By LaSalle, every $\omega$-limit point $X_\infty$
satisfies $\grad\Phi(X_\infty) = 0$.  \textit{Discrete case}: backtracking line
search in \texttt{gravity.py} (lines 659--670) explicitly rejects steps that
increase energy; the force clamp at $F_{\max}=100$ prevents numerical overflow.
\end{proof}

\begin{proof}[Proof of Theorem~\ref{thm:las} (Local Exponential Stability, full)]
Write $z(t) = X(t) - \Xst$ with $\norm{z}$ small.  Expanding:
$\dot z = -\grad\Phi(\Xst + z) = -H^\star z + O(\norm{z}^2).$
Since $H^\star \succ 0$ (Assumption~E2), all eigenvalues of $A^\star = -H^\star$
have real part $-\lambda_i(H^\star) < 0$.  By Lyapunov's indirect method,
$\Xst$ is locally exponentially stable with rate $r = \lmin(H^\star)$.

\textit{Gronwall bound:} decompose $\dot z = -\alpha z + \Delta_{\mathrm{pair}}$
with $\norm{\Delta_{\mathrm{pair}}} \le \Lpair\norm{z}$.  Then:
$\frac{d}{dt}\norm{z}^2 \le -2\alpha\norm{z}^2 + 2\Lpair\norm{z}^2
= -2(\alpha - \Lpair)\norm{z}^2.$
For $\kappa = \alpha - \Lpair > 0$, $\norm{z}^2(t) \le \norm{z}^2(0)e^{-2\kappa t}$.
\end{proof}

\begin{proof}[Proof of Theorem~\ref{thm:phase}(ii) (Growing Mode, full)]
Above $\Cman$, let $u$ be the unit eigenvector of $D^2\Phi_{\mathrm{pair}}(X^+)$
corresponding to $\lambda_u > \alpha$.  Define $W(t) = \tfrac{1}{2}\inner{X(t)-\Xst}{u}^2$.
Differentiating:
$\dot W = \inner{\dot X - 0}{u}\inner{X-\Xst}{u}
= \inner{-\alpha(X-\Xst) - D^2\Phi_{\mathrm{pair}}(X-\Xst)+O(\norm{z}^2)}{u}\inner{X-\Xst}{u}.$
Projecting onto $u$: $\dot W = (-\alpha + \lambda_u)W + O(\norm{z}^3) = cW + O(\norm{z}^3)$
with $c = \lambda_u - \alpha > 0$.  By forward Gronwall:
$W(t) \ge W(0)e^{ct/2}$ for small $t$.  This grows exponentially; no fixed
$\bar\gamma$ prevents it (increasing $\bar\gamma$ shifts $\Cman$ but there always
exist configurations above the shifted $\Cman$ for small enough pairwise distances).
\end{proof}

\begin{proof}[Proof of Proposition~\ref{prop:cs1} (OGD Regret, full)]
Let $\Delta_\gamma = \gamma_{\max} - \gamma_{\min}$.  Standard OGD
\citep{hazan2016oco} with convex $c_t$ and step $\eta_\gamma = \Delta_\gamma/\sqrt{T}$:
$R_T \le \frac{\Delta_\gamma^2}{2\eta_\gamma} + \frac{\eta_\gamma}{2}\sum_t\norm{\nabla_\gamma c_t}^2
\le \frac{\Delta_\gamma\sqrt{T}}{2} + \frac{L_c^2\Delta_\gamma T}{2\sqrt{T}}
= L_c\Delta_\gamma\sqrt{T}.$

\textit{Oracle approximation by power iteration:} maintain unit vector $v_t$;
update $\tilde v_{t+1} = (\grad\Phi_{\mathrm{pair}}(X^t + hv_t) -
\grad\Phi_{\mathrm{pair}}(X^t))/h$, renormalize.  After $K$ steps, the
Rayleigh quotient $\hat\lambda_t$ satisfies
$|\hat\lambda_t - \lmax| \le \lmax(\lambda_2/\lmax)^{2K}$.
Choosing $K = O(\log N)$ gives approximation error $O(N^{-1})$.
\end{proof}

\begin{proof}[Proof of Theorem~\ref{thm:or2} (Optimal Policy, full derivation)]
First-order condition of the Bellman equation w.r.t.\ $\gamma$:
$0 = 2\kappa\gamma + \beta\,\partial_\gamma\E[V(X')\mid X,\gamma].$
Using $\partial_\gamma X' = -\eta\grad E(X)$ and $\partial_{X'} V \approx \grad E(X')/(1-\beta)$
(first-order approximation):
$\partial_\gamma\E[V] \approx -\eta/(1-\beta)\inner{\grad E(X')}{\grad E(X)}.$
At leading order $\grad E(X') \approx \grad E(X)$, giving
$2\kappa\gamma^* \approx \beta\eta\norm{\grad E}^2/(1-\beta)$, so
$\gamma^* \approx \frac{\beta\eta\norm{\grad E}^2/(1-\beta)}{2\kappa}.$
The saturation factor $E/(E+\theta)$ with $\theta = \kappa/\beta$ ensures
$\gamma^*\to 0$ as $E\to 0$ (no intervention needed in the normal regime).
\end{proof}

\begin{proof}[Proof of Lemma (Force-Clamp Regularisation -- Appendix to Proof~7)]
The clamped force $F_{\mathrm{clamp}} = \Pi_{B(F_{\max})}(F_{\mathrm{raw}})$ is
the projection onto the closed ball of radius $F_{\max}$.  Projection onto a
convex set is non-expansive:
$\norm{F_{\mathrm{clamp}}(X) - F_{\mathrm{clamp}}(Y)} \le \norm{F_{\mathrm{raw}}(X) - F_{\mathrm{raw}}(Y)}$.
On any sublevel set $\mathcal{L}_c$, log-repulsion coercivity gives minimum pairwise
distance $\delta_{\min}(c) > 0$, so the Jacobian norm of
$F_{\mathrm{raw}}$ is bounded by $F_{\max}/\delta_{\min}(c) < \infty$.
The contraction condition $\alpha > F_{\max}/\delta_{\min}(c)$ is achievable
for any finite $c$ by increasing $\alpha$.
\end{proof}

\section{Computational Details}
\label{app:compute}

\subsection{GravityEngine Parameter Table}

\begin{center}
\begin{tabular}{llll}
\toprule
Parameter & Symbol & Default value & Role \\
\midrule
Radial spring & $\alpha$ & $0.1$ & Mean-reversion strength \\
Pairwise coupling & $\gamma$ & $1.0$ & Interaction scale \\
Gaussian range & $\sigma$ & $1.0$ & Attraction length-scale \\
Repulsion coefficient & $\lambda$ & $0.1$ & Prevents collapse \\
Softening & $\varepsilon$ & $10^{-5}$ & Regularises log-repulsion \\
Force clamp & $F_{\max}$ & $100$ & Bounds pairwise Jacobian \\
Step size & $\eta$ & $0.02$ & Euler step \\
Backtracking factor & --- & $0.5$ & Armijo line-search halving \\
\bottomrule
\end{tabular}
\end{center}

\subsection{Pairwise Force Kernel}

The erf-based pairwise potential between agents $i$ and $j$ at distance
$r = \norm{x_i - x_j}$ is:
\[
  \phi(r) = \gamma\frac{\sigma\sqrt{\pi}}{2}\erf\!\left(\frac{r}{\sigma}\right)
            - \gamma\lambda\log(r + \varepsilon).
\]
The attraction force $-\phi'(r)\hat r$ is Gaussian: $\phi_1'(r) = \gamma e^{-r^2/\sigma^2}$.
The repulsion force is long-range logarithmic: $\phi_2'(r) = \gamma\lambda/(r+\varepsilon)$.
The net force is attractive at long range and repulsive at short range, giving a
stable cluster configuration at the equilibrium distance
$r^* = \sigma\sqrt{\log(1/\lambda)}$ (implicit solution of $\phi'(r^*) = 0$).

\subsection{Finite-Difference HVP}

The Hessian-vector product $D^2\Phi_{\mathrm{pair}}(X)v$ required by the power
iteration is computed via the symmetric finite difference:
\[
  D^2\Phi_{\mathrm{pair}}(X)v \approx
  \frac{\grad\Phi_{\mathrm{pair}}(X + hv) - \grad\Phi_{\mathrm{pair}}(X - hv)}{2h},
  \quad h = 10^{-4}.
\]
This requires two gradient evaluations ($O(N^2 d)$ each) per power iteration; the
approximation error is $O(h^2\norm{D^3\Phi})$, negligible for $h = 10^{-4}$.

\section{Data Sources and Processing}
\label{app:data}

% [DATA SECTION — to be completed]
See Section~\ref{sec:empirical} and the replication package for full details.

\section{Additional Crisis Window: 2022 Rate Shock}
\label{app:2022}

% [DATA SECTION — to be completed]
Results for the 2022 rate-shock window available in the replication package.

\end{appendices}

%% ====================================================================
%%  Bibliography (placeholder)
%% ====================================================================
\bibliographystyle{plainnat}
% \bibliography{refs}
\begin{thebibliography}{99}

\bibitem{adrian2016covar}
Adrian, T. and Brunnermeier, M.K. (2016).
CoVaR.
\textit{American Economic Review}, 106(7), 1705--1741.

\bibitem{brownlees2017srisk}
Brownlees, C. and Engle, R.F. (2017).
SRISK: A conditional capital shortfall measure of systemic risk.
\textit{Review of Financial Studies}, 30(1), 48--79.

\bibitem{hazan2016oco}
Hazan, E. (2016).
\textit{Introduction to Online Convex Optimization}.
Foundations and Trends in Optimization, 2(3--4), 157--325.

\bibitem{khalil2002nonlinear}
Khalil, H.K. (2002).
\textit{Nonlinear Systems}, 3rd ed.
Prentice Hall.

\bibitem{lasalle1960}
LaSalle, J.P. (1960).
Some extensions of Liapunov's second method.
\textit{IRE Transactions on Circuit Theory}, 7(4), 520--527.

\bibitem{lohmiller1998contraction}
Lohmiller, W. and Slotine, J.-J.E. (1998).
On contraction analysis for non-linear systems.
\textit{Automatica}, 34(6), 683--696.

\bibitem{trefethen1997numerical}
Trefethen, L.N. and Bau, D. (1997).
\textit{Numerical Linear Algebra}.
SIAM.

\bibitem{acemoglu2015systemic}
Acemoglu, D., Ozdaglar, A. and Tahbaz-Salehi, A. (2015).
Systemic risk and stability in financial networks.
\textit{American Economic Review}, 105(2), 564--608.

\bibitem{allen2012systemic}
Allen, L., Bali, T. and Tang, Y. (2012).
Does systemic risk in the financial sector predict future economic downturns?
\textit{Review of Financial Studies}, 25(10), 3000--3036.

\bibitem{bertsekas2012dynamic}
Bertsekas, D.P. (2012).
\textit{Dynamic Programming and Optimal Control}, 4th ed.
Athena Scientific.

\bibitem{brunnermeier2009deciphering}
Brunnermeier, M.K. (2009).
Deciphering the liquidity and credit crunch 2007--2008.
\textit{Journal of Economic Perspectives}, 23(1), 77--100.

\bibitem{elliott2014financial}
Elliott, M., Golub, B. and Jackson, M.O. (2014).
Financial networks and contagion.
\textit{American Economic Review}, 104(10), 3115--3153.

\bibitem{greenwald2020financial}
Greenwald, D.L., Jordà, O., Schularick, M. and Taylor, A.M. (2020).
Financial crisis recoveries.
\textit{AEA Papers and Proceedings}, 110, 124--129.

\bibitem{may2008complex}
May, R.M., Levin, S.A. and Sugihara, G. (2008).
Complex systems: Ecology for bankers.
\textit{Nature}, 451, 893--895.

\bibitem{shalev2012online}
Shalev-Shwartz, S. (2012).
Online learning and online convex optimization.
\textit{Foundations and Trends in Machine Learning}, 4(2), 107--194.

\end{thebibliography}

\end{document}
